\chapter{Literature Survey}\label{ch:lit}

\begin{chapterabstract}
This chapter surveys exact and approximation algorithms for dispersion,
with emphasis on points in convex position, and describes the two papers on
which the project builds: the bounded search tree and dynamic program of
Singireddy, Basappa and Mitchell, and the independent-set approach of
Tkachenko and Wang.
\end{chapterabstract}

\section{Dispersion in general}

For a set of $n$ points with arbitrary distances, $k$-dispersion is NP-hard
even when the distances satisfy the triangle inequality, and for Euclidean
points in the plane it is NP-hard as well~\citep{wangkuo}. Ravi, Rosenkrantz
and Tayi~\citep{ravi} showed that no polynomial algorithm achieves a ratio
better than $1/2$ for metric distances unless $\mathrm{P}=\mathrm{NP}$, and gave a
$1/2$-approximation; the same guarantee had been obtained earlier by
Tamir~\citep{tamir}. The algorithm is the farthest-first greedy rule, which is
also the classical $2$-approximation for $k$-center~\citep{gonzalez}.

For exact algorithms in the plane, Akagi et al.~\citep{akagi} gave an
$n^{O(\sqrt k)}$-time algorithm, together with linear-time algorithms for the
special cases in which the points lie on a line or on a circle; Araki and
Nakano~\citep{arakinakano} improved the line case to $O(\log n)$ time after
preprocessing. Chen and Dumitrescu~\citep{chendumitrescu} recently gave simple
algorithms for small~$k$. Agarwal, Overmars and Sharir~\citep{agarwal} studied
maximally separated sets in the plane, and Fekete and
Meijer~\citep{feketemeijer} considered dispersion with an average-distance
objective.

\section{Dispersion for points in convex position}

For $k=3$, Horiyama et al.~\citep{horiyama} gave an $O(n^2\log n)$ algorithm
for the Euclidean metric, and Kobayashi et al.~\citep{kobayashi} an $O(n^2)$
algorithm for points in convex position. Mishra, Rao and
Das~\citep{mishra} studied an approximation for points on a convex polygon.

\paragraph{Singireddy, Basappa and Mitchell.}
The paper of Singireddy, Basappa and Mitchell~\citep{sbm} is the starting point
of this project. It contains three results.
\begin{enumerate}[leftmargin=*,itemsep=2pt]
\item A \emph{bounded search tree}. The first point is placed at a vertex of the
hull; afterwards the chosen points always form a contiguous arc, and the next
point is the first vertex beyond one of the two ends of this arc that is far
enough from all chosen points. Branching on the end gives at most $2^{k}$
paths per starting vertex, and binary search over the candidate distances
gives an $O(2^kn^2\log^2n)$ algorithm, which is fixed-parameter tractable in
$k$. Its correctness rests on the claim that some starting vertex and some
sequence of ends reach $k$ points whenever a solution exists (Claim~2
of~\citep{sbm}). This claim is the subject of \cref{sec:sbm}.
\item A \emph{dynamic program} over Delaunay triangulations of the solution,
running in $O(n^4k^2)$ time. It was the first polynomial algorithm for
general~$k$.
\item An $O(\log n)$-time $1/(2\sqrt2)$-approximation for $k=3$ that uses only the
extreme points of the polygon.
\end{enumerate}

\paragraph{Tkachenko and Wang.}
Tkachenko and Wang~\citep{tw25} observed that the decision version of
dispersion is a maximum-independent-set problem in the unit-disk graph of the
points, and solved it by dynamic programming over triangulations; for
dispersion their algorithm runs in $O(n^{37/11}\log n)$ time, and for $k=3$
the decision problem can be solved in $O(n\log n)$ time. In a recent
paper~\citep{tw26} they compute a maximum-weight independent set of $n$ disks
in convex position in $O(n^3\log n)$ time, which gives an $O(n^3\log^2 n)$
algorithm for dispersion. A key ingredient~\citep[Lemma~12]{tw26} is that, for
a fixed chord, the constraints imposed by the triangle on the other side of
the chord are monotone along the bisector of the chord; all subproblems of a
chord can therefore be evaluated by sorting and binary search. The dynamic
program of \cref{ch:exact} uses the same monotonicity.

\section{Summary of known bounds}

\Cref{tab:known} summarises the exact algorithms for points in convex position,
together with the results of this project.

\begin{table}[htb]
\centering
\caption{Exact algorithms for $k$-dispersion of $n$ points in convex position.
The general bound of~\citet{tw25} is as reported in~\citet{tw26}.}
\label{tab:known}
\small\setlength{\tabcolsep}{4pt}
\begin{tabular}{@{}lll@{}}
\toprule
Reference & Time & Remarks\\
\midrule
\citet{sbm} & $O(2^k n^2\log^2 n)$ & search tree; correct for $k\le4$ by \cref{thm:sbm4}\\
\citet{sbm} & $O(n^4k^2)$ & Delaunay dynamic program\\
\citet{tw25} & $O(n^{37/11}\log n)$ & via unit-disk independent set\\
\citet{tw26} & $O(n^3\log^2 n)$ & via weighted independent set of disks\\
\citet{kobayashi} & $O(n^2)$ & $k=3$\\
\citet{tw25} & $O(n\log n)$ & $k=3$, decision version\\
\citet{akagi} & $O(n)$ after sorting & points on a circle\\
\midrule
\textbf{This project} (\cref{thm:arcopt}) & $O(n\log^2n)$ expected & unimodal point sets\\
\textbf{This project} (\cref{thm:opt}) & $O(n^3\log k\log n)$ & decision in $O(n^3\log k)$\\
\bottomrule
\end{tabular}
\end{table}

\section{Research gaps}

The survey leaves the following gaps, which this project addresses.
\begin{itemize}[leftmargin=*,itemsep=2pt]
\item The correctness of the fixed-parameter algorithm of~\citet{sbm} depends on
a claim without a complete proof.
\item The fast algorithms for $k=3$ exploit structure tied to the first far
points of a point; it was not known which of these properties survive for
larger~$k$.
\item Apart from points on a circle, no class of point sets in convex
position was known on which dispersion can be solved in near-linear time.
\item The cubic algorithms for general $k$ pay a $\log n$ factor per decision even
when $k$ is small.
\end{itemize}
